% =============================================================
%  ENLACE A: [COMPLETAR]
%  Responsables: [COMPLETAR]
%  Misma estructura que el Enlace B (secciones/03_enlace_B.tex).
%  1) Completar los números en datos/enlace_A.tex
%  2) Reemplazar cada \completar[...] por el texto propio
%  3) Figuras en figuras/A_*.png|jpg
% =============================================================
\renewcommand{\seccorta}{Enlace A}
\section{Enlace A: \EAnombre}

% ---------------------------------------------------------------
\begin{frame}{Definición del enlace y relevamiento geográfico}
  \small
  \begin{tabularx}{\textwidth}{@{}lLL@{}}
    \toprule
     & \textbf{Estación A: \EAunoNombre} & \textbf{Estación B: \EAdosNombre} \\ \midrule
    Ubicación & \EAunoUbic & \EAdosUbic \\
    Latitud / Longitud & \EAunoLat\ / \EAunoLon & \EAdosLat\ / \EAdosLon \\
    Cota del terreno (s.n.m.) & \EAunoCota & \EAdosCota \\
    Altura máxima disponible & \EAunoMax & \EAdosMax \\
    Altura de instalación propuesta & \textbf{\EAunoAlt} & \textbf{\EAdosAlt} \\
    Cota total de la antena & \EAunoAntena & \EAdosAntena \\
    Azimut & \EAazUnoDos\ (A$\to$B) & \EAazDosUno\ (B$\to$A) \\
    Elevación (tilt) & \EAtiltUno & \EAtiltDos \\
    \bottomrule
  \end{tabularx}
  \vspace{6pt}
  \begin{columns}[T]
    \column{0.32\textwidth}
    \begin{block}{Distancia}\centering\EAdist\end{block}
    \column{0.32\textwidth}
    \begin{block}{$\Delta$ terreno}\centering\EAdifTerreno\end{block}
    \column{0.32\textwidth}
    \begin{block}{$\Delta$ antenas}\centering\EAdifAntenas\end{block}
  \end{columns}
\end{frame}

% ---------------------------------------------------------------
\begin{frame}{Perfil del terreno y obstáculos}
  \begin{columns}[c]
    \column{0.52\textwidth}
    \begin{itemize}
      \item Herramientas: \completar[UISP / Google Earth]
      \item Tipo de topografía: \completar
      \item Obstáculo crítico: \EAobstaculo, a $d_1=\EAdUno$ km de la Est.~A
      \item Diferencias entre UISP y Google Earth 3D: \completar
    \end{itemize}
    \column{0.45\textwidth}\centering
    \figura{A_perfil.png}{\linewidth}
  \end{columns}
\end{frame}

% ---------------------------------------------------------------
\begin{frame}{Línea de vista y primera zona de Fresnel}
  \begin{columns}[c]
    \column{0.55\textwidth}
    \small
    \[
      R_1 = 17.32\sqrt{\frac{\EAdUno\times\EAdDos}{\EAfrecNum\times D}} = \textbf{\EArFresnel}
    \]
    \begin{itemize}
      \item Despeje mínimo (60\,\% de $R_1$): \EArSesenta
      \item Despeje medido: \EAdespeje
      \item Despeje logrado: \EAdespejePct\ de $R_1$
    \end{itemize}
    \column{0.42\textwidth}\centering
    \figura{A_google_earth.jpg}{\linewidth}
  \end{columns}
\end{frame}

% ---------------------------------------------------------------
\begin{frame}{Determinación de las alturas mínimas de antena}
  \begin{columns}[c]
    \column{0.55\textwidth}
    \small
    \begin{itemize}
      \item Alturas mínimas: \EAunoAlt\ / \EAdosAlt\ $\Rightarrow$ señal esperada \EAprxSim
      \item Efecto de subir +5 m: \completar
      \item Efecto de subir +10 m: \completar
      \item Conclusión: \completar
    \end{itemize}
    \column{0.42\textwidth}\centering
    \figura{A_senal_uisp.png}{\linewidth}
  \end{columns}
\end{frame}

% ---------------------------------------------------------------
\begin{frame}{Selección de frecuencia y marco regulatorio}
  \begin{block}{Banda elegida: \EAfrec}
    Justificación: \completar
  \end{block}
  \small
  \begin{tabularx}{\textwidth}{@{}lL@{}}
    \toprule
    Banda & \completar \\
    Condición de uso & \completar \\
    Normativa & \completar \\
    Restricciones & \completar \\
    \bottomrule
  \end{tabularx}
\end{frame}

% ---------------------------------------------------------------
\begin{frame}{Selección del ancho de canal}
  \begin{columns}[c]
    \column{0.45\textwidth}\centering
    \begin{tabular}{cc}
      \toprule
      \textbf{Ancho} & \textbf{Capacidad} \\ \midrule
      10 MHz & \completar \\
      20 MHz & \completar \\
      40 MHz & \completar \\
      80 MHz & \completar \\
      \bottomrule
    \end{tabular}
    \column{0.52\textwidth}
    \begin{alertblock}{Selección: \EAcanal}
      \completar[justificación]
    \end{alertblock}
  \end{columns}
\end{frame}

% ---------------------------------------------------------------
\begin{frame}{Selección del equipamiento: \EAequipo}
  \small
  \begin{columns}[T]
    \column{0.52\textwidth}
    \begin{tabularx}{\linewidth}{@{}lL@{}}
      \toprule
      Modelo & \completar \\
      Banda & \completar \\
      Potencia máx. TX & \completar \\
      Sensibilidad & \completar \\
      Modulación / MCS & \completar \\
      Anchos de canal & \completar \\
      Throughput máx. & \completar \\
      Antena / ganancia & \EAganancia \\
      Interfaz / alimentación & \completar \\
      \bottomrule
    \end{tabularx}
    \column{0.45\textwidth}
    \textbf{Justificación}
    \begin{itemize}
      \item Distancia: \completar
      \item Capacidad: \completar
      \item Frecuencia: \completar
      \item Interferencias: \completar
      \item Margen: \completar
      \item Costo/disponibilidad: \completar
    \end{itemize}
  \end{columns}
\end{frame}

% ---------------------------------------------------------------
\begin{frame}{Capacidad requerida}
  \begin{itemize}
    \item Capacidad mínima definida: \EAcapReq
    \item Justificación del tráfico: \completar
    \item Capacidad simulada con ruido: \EAcapRuido
  \end{itemize}
  \begin{block}{Verificación}
    \centering Capacidad disponible (\EAcapRuido) $>$ requerida (\EAcapReq)\,?
  \end{block}
\end{frame}

% ---------------------------------------------------------------
\begin{frame}{Atenuación de espacio libre y presupuesto de RF}
  \begin{columns}[T]
    \column{0.47\textwidth}
    \[
      A_0 = 92.44 + 20\log f + 20\log d = \textbf{\EAfspl}
    \]
    \small Extremo más crítico y por qué: \completar
    \column{0.5\textwidth}\small
    \begin{tabular}{@{}lr@{}}
      \toprule
      $P_{TX}$ & \EAptx \\
      $L_{TX}$, $L_{RX}$ & \completar \\
      $G_{TX}$ = $G_{RX}$ & \EAganancia \\
      $A_0$ & \EAfspl \\
      $A_{obs}$ & \completar \\ \midrule
      $P_{RX}$ teórica & \EAprxTeo \\
      $P_{RX}$ simulada & \EAprxSim \\
      \bottomrule
    \end{tabular}
  \end{columns}
  \vspace{6pt}
  \begin{alertblock}{Discrepancia teoría vs.\ simulación}
    \completar
  \end{alertblock}
\end{frame}

% ---------------------------------------------------------------
\begin{frame}{PIRE, sensibilidad y margen de enlace}
  \[ \text{PIRE} = P_{TX} + G_{TX} = \textbf{\EApire} \]
  Verificación regulatoria: \completar
  \vspace{6pt}\\
  \begin{tabular}{@{}lcc@{}}
    \toprule
    \textbf{MCS} & $S_{RX}$ & Margen \\ \midrule
    Máxima (256QAM) & \completar & \EAmargenMax \\
    Media (16QAM)   & \completar & \EAmargenMed \\
    Básica (BPSK)   & \completar & \completar \\
    \bottomrule
  \end{tabular}
  \vspace{4pt}\\
  Rate adaptation / conclusión: \completar
\end{frame}

% ---------------------------------------------------------------
\begin{frame}{Ruido, SNR, pérdidas adicionales e interferencias}
  \[ \text{SNR} = P_{RX} - P_{ruido} = \EAprxSim - (\EAruido) = \textbf{\EAsnr} \]
  \begin{itemize}
    \item Modulación alcanzable con ese SNR: \EAmcs
    \item Pérdidas por obstáculos ($A_{obs}$): \completar
    \item Análisis de interferencias: \completar
    \item Extremo más crítico: \completar
  \end{itemize}
\end{frame}

% ---------------------------------------------------------------
\begin{frame}{Orientación, montaje, cableado y protección}
  \small
  \begin{columns}[T]
    \column{0.48\textwidth}
    \begin{tabular}{@{}lcc@{}}
      \toprule
       & Azimut & Tilt \\ \midrule
      Est. A & \EAazUnoDos & \EAtiltUno \\
      Est. B & \EAazDosUno & \EAtiltDos \\
      \bottomrule
    \end{tabular}\\[4pt]
    Polarización: \completar\\
    Procedimiento de alineación: \completar
    \column{0.48\textwidth}
    \begin{itemize}
      \item Montaje Est.~A: \completar
      \item Montaje Est.~B: \completar
      \item Cable: tipo \completar, longitud \EAcable
      \item Alimentación: \completar
      \item Protección / PAT: \completar
    \end{itemize}
  \end{columns}
\end{frame}

% ---------------------------------------------------------------
\begin{frame}{Optimización del enlace}
  \begin{block}{Configuración propuesta}
    \EAequipo, \EAunoAlt/\EAdosAlt, canal de \EAcanal
  \end{block}
  \begin{itemize}
    \item Mayor ancho de canal: \completar
    \item Mayor altura: \completar
    \item Antenas de mayor ganancia: \completar
    \item Otra frecuencia: \completar
  \end{itemize}
\end{frame}
